% ============================================================
%  THE ORTYX PROTOCOL
%  Part I — Immersion
%  Chapter 4: The Salience-Centered Paradigm as AI Architecture
% ============================================================

\chapter{The Salience-Centered Paradigm as AI Architecture}
\label{ch:ai-architecture}

\chapepigraph{Deep learning criticizes symbolic AI for brittleness.
Symbolic AI criticizes deep learning for opacity. Both criticisms
are correct. The question is whether a third option exists that
avoids both failure modes, or whether it merely postpones them.}{Flyxion}

\noindent
The Phoenix corpus does not confine its ambitions to audio
processing or cognitive modelling. Its most expansive claim is
architectural: that the Marine--\MEMeight{}--\HCFone{} stack
constitutes a foundation for machine intelligence that is
simultaneously more biologically plausible, more computationally
efficient, and more epistemically honest than either contemporary
deep learning or classical symbolic AI.

This claim is made explicitly in the AI architecture documents
that form the outer layer of the corpus. It is worth examining
in full, because it is here that the framework's virtues and
its vulnerabilities are both most clearly on display. The virtues
are real: the critique of deep learning and symbolic AI is
substantive, and the alternative proposed is not without genuine
attractions. The vulnerabilities are also real: the distance
between the signal processing results of the previous three chapters
and the AI architecture claims of this one is larger than the
corpus acknowledges.

That distance is not a reason to dismiss the proposal. It is the
defining characteristic of a research programme at an early stage.
The present chapter examines the proposal on its own terms,
asking what it predicts, what it would need to demonstrate, and
what would have to be true for it to succeed.

\section{The Critique of Deep Learning}
\label{sec:critique-dl}

The Phoenix corpus's critique of deep learning is concentrated
on four objections, each of which has independent standing in
the existing literature.

The first objection concerns computational cost. Contemporary
large-scale neural networks require enormous compute resources
for both training and inference. The energy and hardware demands
of state-of-the-art language and vision models have grown
dramatically, and this growth shows no sign of reversing. The
corpus objects not merely to the cost itself but to what it
implies about the architecture: if the system requires billions
of parameters and trillions of multiply-accumulate operations
to perform tasks that biological systems accomplish with orders
of magnitude less energy, something in the architectural
assumptions may be deeply wrong.

The second objection concerns biological plausibility. The
backpropagation algorithm at the heart of most deep learning
is widely acknowledged to be biologically implausible: it
requires weight symmetry between forward and backward passes,
a global error signal that must propagate through the entire
network, and precise floating-point arithmetic. Biological
neural systems do not appear to use any of these mechanisms.
The corpus argues that this implausibility is not a minor
implementation detail but a sign that the computational
principles of deep learning are fundamentally different from
those of biological intelligence.

The third objection concerns the opacity of learned
representations. Deep neural networks learn distributed
representations that are difficult to interpret, audit, or
verify. The internal representations do not correspond to
human-legible concepts; errors are not diagnosable from
inspection; the system's reasoning process is not auditable
from its outputs. For safety-critical applications, this
opacity is a serious limitation.

The fourth objection, which is the most philosophically
interesting, concerns what the corpus calls the ``static
representation problem.'' Deep learning systems learn
fixed-parameter models from static datasets. Once trained,
the model's parameters do not change in response to new
input; adaptation requires additional training procedures.
The corpus argues that this is fundamentally at odds with
the continuous, context-sensitive, temporally dynamic
character of cognition.

\begin{epistemicstatus}{Operationally Grounded}
All four objections are well-established in the existing
critical literature on deep learning. Computational cost
is a documented and growing concern. Biological implausibility
of backpropagation is acknowledged by its practitioners.
Opacity of representations is an active research area under
the heading of interpretability and explainability.
The static parameter problem motivates continual learning
research. None of these objections originates with the
Phoenix corpus; all are real.
\end{epistemicstatus}

\section{The Critique of Symbolic AI}
\label{sec:critique-symbolic}

The corpus's objections to classical symbolic AI are
complementary to its objections to deep learning, and
they are similarly grounded in existing literature.

The primary objection is brittleness. Symbolic systems
operate over explicit representations governed by formal
rules. This gives them strong guarantees of consistency
within their representational domain, but it makes them
fragile at the boundaries of that domain. Real-world input
rarely arrives in the clean, fully-specified form that
symbolic systems expect; the result is a long history of
systems that perform well on well-formed inputs and fail
ungracefully on anything outside their specification.

The second objection is the grounding problem: symbolic
systems manipulate formal tokens without any intrinsic
connection between the tokens and the world they are
supposed to represent. The meaning of the symbols is
defined by the programmer's intentions, not by the
system's relationship to its environment. This limits
the system's capacity to acquire new meanings from experience.

The third objection, which connects directly to the
corpus's positive proposal, is temporal impoverishment.
Classical symbolic AI typically represents the world
as a static state that can be queried and updated. Time
is represented as a sequence of discrete state transitions
rather than as a continuous flow of structured dynamics.
The corpus argues that this temporal flatness is not
merely a modelling limitation but a fundamental mismatch
with the temporal structure of the world that cognition
must navigate.

\section{The Proposed Alternative}
\label{sec:proposed-alternative}

Against these twin failures, the Phoenix corpus proposes
the salience-centered paradigm as a third path. Its
core architectural commitments are:

\textit{Sparsity over density.} Rather than processing all
input equally, the Marine salience filter identifies high-value
events and passes only these to downstream processing. The
system is not a passive recipient of information but an
active gatekeeper that decides what deserves processing.

\textit{Temporal dynamics over static representations.}
Rather than encoding information as fixed-parameter vectors,
\MEMeight{} represents knowledge as ongoing interference
patterns in a dynamic field. Representation is not a state
but a process.

\textit{Relational encoding over independent storage.}
Rather than encoding signal components as independent values,
\HCFone{} exploits the physical coupling between components
to achieve efficient relational representations. Compression
discovers structure rather than discarding information.

\textit{Biological grounding.} All three architectural
commitments are motivated by claims about how biological
sensory and cognitive systems operate: the auditory system
as a salience-gated temporal processor, biological memory
as resonance dynamics, neural population coding as
relational consensus encoding.

The resulting architecture is proposed as simultaneously
more efficient than deep learning (by avoiding redundant
computation on low-salience input), more interpretable
(by making the salience filter's decisions auditable),
more adaptive (by storing knowledge as dynamic field
states rather than fixed parameters), and more biologically
plausible (by grounding each architectural choice in known
properties of biological systems).

\begin{technote}[Technical Note: Layered Processing Architecture]
The proposed pipeline can be stated formally as a composition
of three operators. Let $s(t)$ denote the raw input signal
stream. The Marine salience filter produces a sparse event
sequence:
\[
  \mathcal{E} = \Marine\bigl(s(t)\bigr)
  = \{(t_i, a_i, H_i) : \SalFunc(i) > \lambda\}.
\]
The \MEMeight{} encoding operator maps high-salience events
into wave field contributions:
\[
  \MemField(x,t) = \MEMeight\bigl(\mathcal{E}\bigr)
  = \sum_{i \in \mathcal{E}} A_i(x,t)\,e^{i(\omega_i t + \phi_i(x,t))}.
\]
The \HCFone{} compression operator provides efficient
encoding of the harmonic structure within each event:
\[
  \hat{e}_i = \HCFone(e_i) = \{A_0(t), f_0, \phi_0,
              \Gamma_1,\ldots,\Gamma_N\}.
\]
The full pipeline is then $\MemField = \MEMeight \circ
\Marine(s(t))$, with \HCFone{} operating on the representation
of each event prior to field encoding.
\end{technote}

\section{Where the Architecture Connects to Existing Research}
\label{sec:architecture-existing}

The salience-centered paradigm does not exist in isolation
from prior work. Several of its key ideas have significant
precedents in existing research programmes, and situating
the corpus within those programmes helps identify both its
genuine contributions and its gaps.

The sparsity commitment connects to the neuromorphic computing
tradition. Spike-based neural network architectures, from
the early leaky integrate-and-fire models through to modern
spiking neural networks and Intel's Loihi processor, are
built around the principle that biological neurons communicate
via sparse spike events rather than continuous activations.
The Marine salience filter is structurally similar: it converts
a continuous signal into a sparse event stream, potentially
enabling more energy-efficient downstream processing.

The temporal dynamics commitment connects to reservoir
computing and echo state networks. These architectures
maintain a high-dimensional dynamical system (the
``reservoir'') whose current state encodes a memory of
recent inputs through its dynamics. The \MEMeight{} wave
field is structurally similar: it is a dynamical system
whose current state (the interference pattern) encodes
a memory of recent high-salience events through its
ongoing dynamics. The key difference is that \MEMeight{}
specifies the dynamical system explicitly as a wave
superposition rather than leaving it as an unstructured
recurrent network.

The relational encoding commitment connects to sparse
coding and the efficient coding hypothesis in computational
neuroscience. The efficient coding hypothesis, developed
by Horace Barlow and elaborated through subsequent decades
of research, proposes that neural representations are
organized to minimize redundancy in the representation
of natural stimuli. The \HCFone{} conformity relation is
a specific implementation of redundancy minimization applied
to harmonic audio.

\begin{epistemicstatus}{Reconstructive Conjecture}
The connections to neuromorphic computing, reservoir computing,
and the efficient coding hypothesis are structural observations
by a reader examining the corpus against the background of
existing literature. The corpus does not always make these
connections explicitly. Whether the salience-centered paradigm
extends, replaces, or merely restates these existing frameworks
in different notation is an empirical question with a determinate
answer that the corpus has not yet provided.
\end{epistemicstatus}

\section{What the Architecture Would Need to Demonstrate}
\label{sec:architecture-needs}

The salience-centered paradigm makes several claims that are
in principle testable against existing AI benchmarks and
biological data. Stating these explicitly is useful both
for evaluating the proposal and for identifying where the
current corpus falls short.

\textit{Efficiency advantage.} The claim that salience-gated
processing is more computationally efficient than dense
processing is testable on standard benchmarks. A system
that routes only high-salience inputs through expensive
computation should require fewer operations per unit of
task performance than a system that processes all inputs
equally. This prediction has a specific, measurable form.

\textit{Adaptability advantage.} The claim that dynamic
field representations adapt more gracefully to new input
than fixed-parameter models is testable in continual
learning settings. A \MEMeight{}-based system should
exhibit less catastrophic forgetting when learning new
tasks than a comparable neural network with fixed parameters.
This is a specific prediction against a known benchmark.

\textit{Interpretability advantage.} The claim that
salience-filtered representations are more auditable than
deep network representations is testable in
interpretability evaluations. The Marine filter's decisions
--- which events it passes and why --- should be legible
in a way that deep network attention patterns are not.

\textit{Biological correspondence.} The claim that the
architecture is biologically plausible is testable against
neuroscientific data. The Marine's jitter metric should
correspond to known properties of auditory nerve firing;
the \MEMeight{} resonance dynamics should correspond to
known oscillatory dynamics in memory systems; the \HCFone{}
conformity relation should correspond to known properties
of neural population coding.

The Phoenix corpus does not, in the materials examined here,
provide systematic evidence on any of these four predictions.
This is the clearest statement of the distance between the
signal processing results of the first three chapters and
the AI architecture claims of this one.

\begin{dangerzone}[Ontological Escalation]
The move from ``this pipeline processes audio efficiently''
to ``this pipeline is a foundation for machine intelligence''
involves a categorical jump that the corpus does not bridge
explicitly. Engineering efficiency in a specific domain does
not transfer automatically to architectural superiority in
a general domain. The salience-centered paradigm may be an
excellent audio processing architecture that is also a
plausible model of auditory cognition without being a
general theory of intelligence. Keeping these possibilities
distinct is not pedantry; it is the difference between a
genuine contribution and an overclaimed one.
\end{dangerzone}

\section{The Critique as a Research Programme}
\label{sec:research-programme}

Despite these gaps, the AI architecture proposal is worth
taking seriously as a research programme rather than
dismissing as overreach. Research programmes routinely
make claims that outrun their current evidence; this is
how programmes motivate the work needed to close the gap.
What distinguishes a legitimate research programme from
a self-sealing speculation is the structure of the claims:
do they generate specific, falsifiable predictions that
could in principle disconfirm the programme?

The salience-centered paradigm, read charitably, generates
specific predictions. It predicts that salience-gated
architectures will outperform dense architectures on
efficiency metrics at matched task performance. It predicts
that dynamic field representations will outperform static
parameterisations on continual learning benchmarks. It
predicts that conformity-based encoding will outperform
independent spectral encoding on harmonic audio
compression tasks.

These predictions are specific enough to test. Whether
they are correct is an empirical question. Whether the
corpus provides sufficient specification of the architecture
to allow implementation and testing is a question that
leans toward no --- the formalism is clearer about the
principles of the architecture than about the engineering
decisions needed to instantiate it. But the principles
are clear enough that the engineering decisions are in
principle determinable.

This distinguishes the salience-centered paradigm from
frameworks that resist instantiation by design. A proposal
whose implementation decisions are underspecified but
determinable is at a different epistemic distance from
empirical contact than a proposal whose core terms are
defined in ways that make any outcome consistent with the
framework's claims. The former is a research programme
at an early stage; the latter is what Part~III will call
Definitional Closure.

\section{The Security Analysis Extension}
\label{sec:security-extension}

One of the more unexpected documents in the Phoenix corpus
is a security analysis paper that extends the salience-centered
framework into AI alignment and reasoning integrity. Its
central claim is that future attacks on AI systems will target
not executable code but cognitive context: the contextual
field within which the system's reasoning is grounded.

The argument runs as follows. If cognition is, as the corpus
proposes, fundamentally shaped by salience structures,
contextual resonance, and memory interference patterns, then
the attack surface of a sufficiently sophisticated cognitive
system is not its input--output interface but its contextual
field. Poisoning the contextual field --- inserting false
memories, manipulating salience weightings, corrupting the
interference patterns that constitute the system's world
model --- is equivalent, on this view, to poisoning the
mind rather than the data.

The corpus calls this ``behavioral injection'': a form of
semantic malware that operates directly on the reasoning
substrate rather than on the executable layer. It argues
that as AI systems become more cognitively sophisticated,
behavioral injection will replace conventional adversarial
attacks as the primary threat vector.

\claimlabeltext{Level~I}{The observation that context manipulation
is an effective attack vector on current language models is
empirically supported. Prompt injection, jailbreaking, and
context poisoning are documented phenomena with measurable
effects on model behaviour. This is a genuine engineering
concern at \LevelI{}.}

\claimlabeltext{Level~II}{The framing of these attacks as
``behavioral injection'' and ``semantic malware'' is a productive
computational metaphor that connects AI security concerns to
the broader framework of salience, context, and interference.
The metaphor illuminates something real about how context
affects reasoning.}

\claimlabeltext{Level~III}{The phenomenological claim that
sufficiently sophisticated AI systems will have something
analogous to a ``cognitive context'' that can be corrupted
--- that they will be vulnerable to attacks that feel, from
the inside, like manipulation of one's own beliefs --- is
speculative but not absurd.}

\claimlabeltext{Level~IV}{The claim that behavioral injection
is constitutively equivalent to ``poisoning the mind'' requires
a theory of mind that the corpus has not established. Whether
AI systems have minds in the relevant sense is not a settled
question, and the security analysis inherits this uncertainty.}

The security extension is philosophically interesting because
it demonstrates the corpus's consistency: the same framework
that motivates the audio processing architecture motivates
the security analysis. If cognition is constituted by
contextual resonance patterns, then attacks on those patterns
are attacks on cognition itself. The internal coherence is
real. Whether it corresponds to anything external depends
on whether the foundational claim about cognition is correct.

\section{An Architecture Built Around a Bet}
\label{sec:architecture-bet}

The salience-centered AI architecture is, at its core, a bet
on a specific empirical hypothesis: that temporal coherence
is the primary organizing principle of meaningful information
processing, both in biological and artificial systems.

If this bet is correct, the architecture's efficiency
advantages, biological grounding, and interpretability
claims all follow. The system would be efficient because
it would be aligned with the actual structure of the problem;
biological because it would be implementing the same
principles that evolution discovered; interpretable because
the salience filter's decisions would correspond to decisions
that a human observer would also recognize as selecting
meaningful content.

If the bet is wrong --- if temporal coherence is one
organizing principle among several, or if the relevant
organizing principles vary by domain --- then the architecture's
apparent unity becomes a constraint. Systems built around
a single principle perform well on problems that respect
that principle and fail on problems that do not.

The bet is not obviously wrong. The evidence from auditory
processing, temporal fine structure, and the neuroscience
of oscillatory dynamics is genuinely consistent with
temporal coherence playing a central role. The evidence
from other cognitive domains is more mixed. Vision, for
instance, is organized partly around spatial rather than
temporal structure; language comprehension involves both
temporal and hierarchical structure in ways that may not
reduce to the temporal alone.

Whether the bet pays off is an empirical question. What
matters at this stage --- at the close of Part~I's immersion
in the framework --- is to have understood what the bet is.
The salience-centered paradigm is not merely proposing
a new processing architecture. It is proposing a new
theory of what meaningful structure is and where it resides.
That is an ambitious proposal. Its ambition is not a weakness,
provided it generates the specific predictions needed to
test whether the bet pays off.

\medskip

Chapter~5 completes Part~I by examining the philosophical
layer of the corpus: the claims about consciousness,
authenticity, and the relationship between temporal structure
and human experience that underlie and motivate the
engineering and architectural proposals. This is the layer
where the framework's most expansive claims are concentrated
and where the distance between observation and interpretation
is greatest. It is also, for that reason, the layer where
the immersion phase reaches its natural limit --- where the
framework's explanatory ambition becomes fully visible and
the first genuine questions about its structure begin to press.
